\textbf{(i)} A union of prime ideals contains a product exactly when either factor belongs to the union. Its complement is therefore saturated.

Conversely, suppose $S$ is saturated and $x\notin S$. Then $(x)\cap S=\varnothing$, since $ax\in S$ would imply $x\in S$. The prime ideal theorem gives a prime ideal containing $(x)$ and disjoint from $S$. Taking the union over $x\notin S$ gives $A\setminus S$.

\textbf{(ii)} Let $\overline S$ be the complement of the union of primes disjoint from $S$. By (i), it is saturated and contains $S$. The complement of any saturated set containing $S$ is a union of primes disjoint from $S$, so that set contains $\overline S$. This proves minimality and uniqueness. Moreover, $x\in\overline S$ if and only if $(x)$ meets $S$, or equivalently $ax\in S$ for some $a\in A$.

For $S=1+\mathfrak a$, this criterion says $ax\equiv1\pmod{\mathfrak a}$ for some $a\in A$. Thus $\overline S$ is the inverse image of $(A/\mathfrak a)^\times$.
